\documentclass[11pt, letterpaper]{article}
% ----- PAQUETES -----
\usepackage[utf8]{inputenc}
\usepackage[T1]{fontenc}
\usepackage[spanish]{babel}
\usepackage{graphicx}
\usepackage[margin=2.3cm]{geometry}
\usepackage{xcolor}
\usepackage{titlesec}
\usepackage{booktabs}
\usepackage{enumitem}
\usepackage{xurl}
\usepackage{hyperref}
\usepackage{fancyhdr}
\usepackage{amsmath}
\usepackage{amssymb}
\usepackage{array}

% ----- COLORES -----
\definecolor{UnalGreen}{RGB}{0,153,51}
\definecolor{DarkGray}{RGB}{60,60,60}

% ----- SECCIONES -----
\titleformat{\section}
{\normalfont\large\bfseries\color{UnalGreen}}
{\thesection}{1em}{}[{\titlerule[0.8pt]}]
\titleformat{\subsection}
{\normalfont\normalsize\bfseries\color{DarkGray}}
{\thesubsection}{1em}{}
\titlespacing*{\section}{0pt}{8pt}{4pt}
\titlespacing*{\subsection}{0pt}{6pt}{2pt}

% ----- ENLACES -----
\hypersetup{
	colorlinks=true,
	linkcolor=UnalGreen,
	urlcolor=UnalGreen
}

% ----- ENCABEZADO -----
\pagestyle{fancy}
\fancyhf{}
\lhead{\footnotesize Mecánica y Ondas para Biociencias (1000023) -- Laboratorio}
\rhead{\footnotesize Semestre 2026-2}
\cfoot{\thepage}

% ----- LISTAS COMPACTAS -----
\setlist[itemize]{leftmargin=*, itemsep=2pt, topsep=2pt, parsep=0pt}
\setlist[enumerate]{leftmargin=*, itemsep=3pt, topsep=3pt, parsep=0pt}

% ----- CAJA DE NOTA -----
\newcommand{\notabox}[1]{%
	\begin{center}
		\fcolorbox{UnalGreen}{UnalGreen!8}{%
			\parbox{0.94\textwidth}{\small #1}%
		}
	\end{center}
}

\begin{document}
	
	% =========================================================
	% ENCABEZADO
	% =========================================================
	\begin{center}
		\vspace*{-0.8cm}
		\IfFileExists{imagenes/universidad-nacional-de-colombia-sede-bogota-seeklogo.png}{%
			\includegraphics[width=0.15\textwidth]
			{imagenes/universidad-nacional-de-colombia-sede-bogota-seeklogo.png}\\[0.1cm]
		}{}
		
		{\Large\bfseries La Bitácora de Laboratorio}\\[0.1cm]
		{\large\bfseries Qué debe contener y cómo entregarla}\\[0.1cm]
		{\normalsize Mecánica y Ondas para Biociencias (1000023)}
	\end{center}
	
	\notabox{Las bitácoras valen el \textbf{50\,\% de la nota} del laboratorio.
		Son \textbf{grupales}: se entrega una sola por grupo. Se envían por
		\textbf{Google Classroom en formato PDF}, antes de la fecha límite publicada
		para cada práctica.}
	
	% =========================================================
	\section{¿Qué es una bitácora?}
	% =========================================================
	
	La bitácora es el registro de lo que \emph{realmente} hicieron y midieron en el
	laboratorio: qué montaron, qué datos obtuvieron, cómo los procesaron y qué
	concluyen. No es un informe formal: \textbf{no lleva resumen ni marco teórico}.
	La teoría ya está en la guía; la bitácora muestra su trabajo experimental.
	
	\vspace{0.15cm}
	\noindent
	Recomiendo que sea \textbf{concisa, pero no por ello incompleta}. Una bitácora
	corta con todos sus datos, unidades e incertidumbres vale más que una larga con
	texto de relleno.
	
	% =========================================================
	\section{Formato}
	% =========================================================
	
	Pueden elaborarla en el medio que prefieran:
	
	\begin{itemize}
		\item \textbf{Cuaderno físico} (cuadriculado): recomendado para tomar datos
		en el laboratorio y hacer las gráficas preliminares a mano.
		\item \textbf{Digital}: Word, Google Docs, \LaTeX, o la que prefieran.
	\end{itemize}
	
	\noindent
	Sea cual sea el formato, \textbf{lo que se entrega es un único archivo PDF}.
	Si trabajan en el cuaderno, escaneen las páginas con una aplicación de
	escáner y únanlas en un solo PDF. Las fotos deben ser legibles, derechas y con buena
	luz: \textbf{lo que no se puede leer, no se califica}.
	
	% =========================================================
	\section{Estructura de la bitácora}
	% =========================================================
	
	\subsection{Encabezado}
	
	\begin{itemize}
		\item Número y nombre de la práctica (por ejemplo: \emph{Bitácora 1 --
			Plano Inclinado}).
		\item Número de grupo y nombres completos de \textbf{todos} los integrantes.
		\item Fecha de realización de la práctica.
	\end{itemize}
	
	\subsection{Montaje}
	
	Un esquema o dibujo del montaje experimental, con las magnitudes que midieron
	señaladas sobre él ($h$, $L$, $x$, etc.). Puede ser a mano o una foto del
	montaje real, pero debe estar \textbf{rotulado} (es decir señalado qué es qué). Agreguen en pocas líneas qué
	hicieron, sobre todo si se apartaron de lo que dice la guía.
	
	\subsection{Datos experimentales}
	
	Es la parte más importante. Todos los datos que tomaron, organizados en tablas:
	
	\begin{itemize}
		\item Cada columna con el nombre de la magnitud y su \textbf{unidad}:
		$h$ (cm), $t$ (s).
		\item Cada medida con su \textbf{incertidumbre}, e indicando de dónde sale
		(resolución del instrumento, dispersión de medidas repetidas, o ambas).
		\item Los datos \textbf{crudos}, tal como los midieron, aunque alguno les
		parezca ``malo''. No se borran datos: si descartan uno, lo tachan y
		explican por qué.
	\end{itemize}
	
	\notabox{\textbf{Sin unidades o sin incertidumbres, la nota se reduce.} Una
		medida sin su incertidumbre no se puede comparar con nada; un número sin unidad
		no significa nada.}
	
	\subsection{Procesamiento}
	
	Los cálculos que hicieron a partir de los datos: promedios, desviaciones
	estándar, propagación de incertidumbres, variables linealizadas.
	
	\begin{itemize}
		\item Muestren \textbf{al menos un ejemplo de cada cálculo} completo, con la
		fórmula, los valores sustituidos y el resultado. Los demás pueden ir
		directamente en una tabla (no llenen la bitacora con calculos repetitivos).
		\item Reporten los resultados con el \textbf{número correcto de cifras
			significativas}: incertidumbre con una cifra significativa y el valor
		redondeado a la misma posición decimal.
	\end{itemize}
	
	\subsection{Gráficas}
	
	\begin{itemize}
		\item Ejes rotulados con la magnitud y la unidad.
		\item Escalas con valores redondos y que aprovechen el espacio disponible.
		\item Puntos con sus \textbf{barras de error}.
		\item Si ajustan una recta, que se vea sobre la gráfica, junto con su
		pendiente e intercepto (con unidades).
		\item Cada gráfica numerada y con un título o pie breve que diga qué muestra.
	\end{itemize}
	
	\noindent
	Pueden hacerlas a mano en papel milimetrado o con software (Excel, Tracker,
	Python). Las gráficas hechas a mano durante la sesión, para verificar
	tendencias, también pueden incluirse.
	
	\subsection{Análisis de resultados}
	
	Respondan las preguntas de análisis de la guía, \textbf{numeradas} igual que en
	ella. Comparen sus resultados con lo que predice la teoría (por ejemplo, con un
	porcentaje de discrepancia) y discutan si la diferencia es compatible con sus
	incertidumbres.
	
	\subsection{Conclusiones}
	
	Pocas y concretas. Una buena conclusión:
	
	\begin{itemize}
		\item Responde a los \textbf{objetivos} de la práctica.
		\item Se apoya en \textbf{números}: ``obtuvimos $g = (9.6 \pm 0.3)$ m/s$^2$,
		compatible con el valor esperado'' dice mucho más que ``el experimento
		salió bien''.
		\item Identifica fuentes de error concretas y propone cómo mejorar el
		experimento.
	\end{itemize}
	

	% =========================================================
	\section{Entrega}
	% =========================================================
	
	\begin{itemize}
		\item Un solo archivo \textbf{PDF} por grupo, subido a Google Classroom por
		\emph{un} integrante.
		\item Nombre del archivo: \texttt{Bitacora1\_Grupo3.pdf} (con el número de
		bitácora y de grupo correspondientes).
		\item Páginas numeradas.
	\end{itemize}
	
	% =========================================================
	\section{Lista de verificación antes de enviar}
	% =========================================================
	
	\begin{itemize}
		\item[$\square$] ¿Están los nombres de todos los integrantes y el número de
		grupo?
		\item[$\square$] ¿El montaje tiene un esquema rotulado?
		\item[$\square$] ¿Todas las tablas tienen unidades en los encabezados?
		\item[$\square$] ¿Todas las medidas y resultados tienen su incertidumbre?
		\item[$\square$] ¿Las incertidumbres tienen una sola cifra significativa?
		\item[$\square$] ¿Hay al menos un ejemplo completo de cada cálculo?
		\item[$\square$] ¿Las gráficas tienen ejes rotulados, unidades y barras de
		error?
		\item[$\square$] ¿Respondieron todas las preguntas de análisis de la guía?
		\item[$\square$] ¿Es un solo PDF, legible, con el nombre correcto?
	\end{itemize}
	
	
	
\end{document}